\clearpage
\VerdierPhysicalPageAnchor{0153}%
\begingroup
\small
\setlength{\parskip}{0pt}%
\setlength{\emergencystretch}{1.5em}%

\VerdierSemanticAnchor{ch3-3-1-3}
\noindent\textbf{3.1.3.}\quad
Supposons que la catégorie $\mathcal D$ soit une $\mathbb Z$-catégorie
stricte. Les foncteurs de translation de $\mathcal D$,
$X\longmapsto X[n]$, $n\in\mathbb Z$, s'étendent en automorphismes de
$\mathcal{F}\ell_1(\mathcal D)$ :

\VerdierSemanticAnchor{ch3-3-1-3-translation-extension}
\[
\bigl(X\xrightarrow{f}Y\bigr)
\longmapsto
\bigl(X[n]\xrightarrow{f[n]}Y[n]\bigr).
\]
Il résulte immédiatement de la propriété universelle de
$\mathrm A(\mathcal D)$ (3.1.1.1) que ces automorphismes passent au quotient
et définissent sur $\mathrm A(\mathcal D)$ une structure de
$\mathbb Z$-catégorie stricte. Le foncteur
$\mathcal H\colon\mathcal D\to\mathrm A(\mathcal D)$ est un
$\mathbb Z$-foncteur strict.

\VerdierSemanticAnchor{ch3-3-1-4}
\medskip
\noindent\textbf{3.1.4.}\quad
Soient $\mathcal D$ et $\mathcal D'$ deux $\mathbb Z$-catégories strictes
additives et $F\colon\mathcal D\to\mathcal D'$ un $\mathbb Z$-foncteur
additif (tordu par le cocycle trivial). Soit $\mathcal{F}\ell_1(F)$
l'extension de $F$ aux catégories de morphismes :
\[
\mathcal{F}\ell_1(F)\bigl(X\xrightarrow{f}Y\bigr)
=\bigl(FX\xrightarrow{Ff}FY\bigr).
\]
Le foncteur $\mathcal{F}\ell_1(F)$ transforme les morphismes négligeables de
$\mathcal{F}\ell_1(\mathcal D)$ en morphismes négligeables de
$\mathcal{F}\ell_1(\mathcal D')$. Il définit par suite un foncteur

\VerdierSemanticAnchor{ch3-3-1-4-1}
\[
\mathrm A(F)\colon\mathrm A(\mathcal D)\longrightarrow\mathrm A(\mathcal D').
\tag{3.1.4.1}
\]
Le foncteur $\mathrm A(F)$ est additif. Il résulte immédiatement des
propriétés universelles de $\mathrm A(\mathcal D)$ et
$\mathrm A(\mathcal D')$ qu'il est muni canoniquement d'une structure de
$\mathbb Z$-foncteur.

\VerdierSemanticAnchor{ch3-3-1-4-composition}
Soit $F'\colon\mathcal D'\to\mathcal D''$ un $\mathbb Z$-foncteur additif
entre $\mathbb Z$-catégories strictes additives. On a alors

\VerdierSemanticAnchor{ch3-3-1-4-2}
\[
\mathrm A(F')\mathrm A(F)=\mathrm A(F'F).
\tag{3.1.4.2}
\]

\VerdierSemanticAnchor{ch3-3-1-4-3}
Le diagramme ci-après est commutatif :
\[
\begin{tikzcd}[column sep=5em,row sep=2.5em]
\mathcal D \arrow[r,"F"] \arrow[d,"\mathcal H"'] &
\mathcal D' \arrow[d,"\mathcal H'"] \\
\mathrm A(\mathcal D) \arrow[r,"\mathrm A(F)"'] &
\mathrm A(\mathcal D')
\end{tikzcd}
\tag{3.1.4.3}
\]
(les foncteurs $\mathcal H$ et $\mathcal H'$ sont définis en (3.1.2)).

\VerdierSemanticAnchor{ch3-3-1-4-natural-transformation}
Soit de même $m\colon F\to F''$ un morphisme de $\mathbb Z$-foncteurs
additifs entre $\mathcal D$ et $\mathcal D'$. On en déduit, par les mêmes
arguments que précédemment, un morphisme de $\mathbb Z$-foncteurs

\VerdierSemanticAnchor{ch3-3-1-4-4}
\[
\mathrm A(m)\colon\mathrm A(F)\longrightarrow\mathrm A(F'').
\tag{3.1.4.4}
\]
Si $m$ et $m'$ sont deux morphismes composables de $\mathbb Z$-foncteurs
additifs, on a

\VerdierSemanticAnchor{ch3-3-1-4-5}
\[
\mathrm A(m')\mathrm A(m)=\mathrm A(m'm).
\tag{3.1.4.5}
\]
\endgroup
